\chapter{Melindungi Sistem Saraf}\label{ch:g8:protecting-the-nervous-system}

Dua bab sudah memperlihatkan \omterm{def:g5:brain-in-command:system}{sistem sarafnya} itu apa: kabel
yang simpangnya belajar, memutuskan --- dan dapat diperdaya
(\cref{rem:g8:nervous-communication:vulnerable}). Penutup
tahun ini mengubah anatomi itu menjadi perlindungan. Apa
sebenarnya yang dilakukan kurang \omterm{prop:g1:healthy-habits:needs}{tidur}, kebisingan, layar
dan zat psikoaktifnya, simpang demi simpang, terhadap mesin
yang adalah dirimu --- dan seperti apa melindunginya dalam
praktik, pada usia empat belas?

\section{Perawatan: tidur, dan batas indranya}

\begin{proposition}[Tidur adalah perawatan sinapsis]\label{prop:g8:protecting-the-nervous-system:sleep}
Pengarsipan-dan-perbaikan pada
\cref{prop:g5:brain-in-command:protect} punya bacaan
setingkat simpang: dalam \omterm{prop:g1:healthy-habits:needs}{tidur}, \omterm{def:g8:nervous-communication:synapse}{sinapsis} yang menguat pada
hari itu dipadatkan --- latihan yang dibuat menetap
(\cref{ex:g8:nervous-communication:juggler}) sebagian besar
terjadi semalaman --- dan tata graha \omterm{def:g5:brain-in-command:system}{otaknya} membersihkan
keausan kimiawi hari itu. Malam yang pendek meninggalkan
pelajaran kemarin setengah terarsip dan penimbangan besok
yang lamban: \omterm{def:g5:brain-in-command:system}{otak} \omterm{def:g3:stages-of-life:stages}{remajanya}, yang sedang dirangkai ulang
(\cref{rem:g8:puberty:moods}) dan jamnya hanyut makin malam,
paling butuh jam \omterm{prop:g1:healthy-habits:needs}{tidurnya} justru ketika sekolahnya paling pagi.
\end{proposition}
\admitted

\begin{example}[Paradoks belajar ulangnya]\label{ex:g8:protecting-the-nervous-system:revision}
Dua murid menjelang sebuah ujian: yang satu belajar ulang
sampai tengah malam lalu terus sampai pukul dua; yang lain
belajar ulang sampai pukul sepuluh lalu \omterm{prop:g1:healthy-habits:needs}{tidur}. Yang kedua
biasanya menang --- bahannya menyeberangi simpangnya
\emph{lalu dipadatkan}; jam tambahan yang pertama justru
mengorbankan pengarsipan yang membuat belajar ulangnya
melekat. ``Tidurkanlah dahulu'' adalah rekayasa \omterm{def:g8:nervous-communication:synapse}{sinapsis}.
\end{example}

\begin{proposition}[Keras adalah sebuah takaran]\label{prop:g8:protecting-the-nervous-system:noise}
\omterm{def:g5:first-look-at-cells:cell}{Sel} indera \omterm{def:g1:the-five-senses:sense}{telinganya} --- penerima sehalus rambut yang
mengubah bunyi menjadi \omterm{def:g8:nervous-communication:neuron}{pesan saraf} --- tidak tumbuh kembali
(\cref{prop:g1:the-five-senses:fragile} mengatakannya dengan
lembut; inilah mekanismenya). Bunyi yang keras membunuhnya
menurut takaran: kenyaringan dikali lama paparannya. Konser
dan penyuara \omterm{def:g1:the-five-senses:sense}{telinga} pada kenyaringan penuh membelanjakan,
jam demi jam, sebuah modal penerima yang tidak dipulihkan
\omterm{prop:g1:healthy-habits:needs}{tidur} mana pun --- desingan sesudah sebuah malam yang keras
(\cref{ex:g1:the-five-senses:loud}) adalah peringatan
takaran berlebihnya. Aturannya adalah sebuah anggaran: lebih
rendah, lebih pendek, dengan jeda --- dan sumbat \omterm{def:g1:the-five-senses:sense}{telinga} di
tempat pengeras suaranya berkuasa.
\end{proposition}
\admitted

\section{Zat yang memperdaya simpangnya}

\begin{proposition}[Zat psikoaktif bekerja di sinapsis]\label{prop:g8:protecting-the-nervous-system:substances}
Setiap zat yang mengubah suasana hati, penginderaan atau
kesiagaan bekerja di celah pada
\cref{def:g8:nervous-communication:synapse} --- menirukan
sebuah pengantar pesan, menghadang salah satunya, atau
memaksakan pelepasannya:
\begin{itemize}
  \item \emph{alkohol} meredam simpang seluruh sistemnya:
        gelungnya melambat dari ujung ke ujung ---
        pertimbangan, \omterm{def:g5:brain-in-command:reaction}{waktu reaksi} (meter pada
        \cref{ex:g5:brain-in-command:driver} memanjang),
        keseimbangan, ingatan; simpang \omterm{def:g5:brain-in-command:system}{otak} \omterm{def:g3:stages-of-life:stages}{remajanya} yang
        masih dibangun menerima peredaman itu paling dalam;
  \item \emph{nikotin tembakaunya} menirukan sebuah
        pengantar pesan alami dan, bila diulang, merangkai
        ulang simpang ganjarannya menuju \emph{\omterm{prop:g5:healthy-choices:substances}{kecanduan}}
        (jebakan pada \cref{prop:g5:healthy-choices:substances},
        kini dengan mekanismenya: sinapsisnya menyesuaikan
        diri untuk mengharapkan takarannya --- lalu
        memprotes tanpanya);
  \item \emph{ganja} menirukan keluarga pengantar pesan yang
        lain, mengeruhkan persis simpang ingatan, dorongan
        dan pengaturan waktunya --- sekali lagi paling mahal
        pada \omterm{def:g5:brain-in-command:system}{otak} yang sedang dibangun;
  \item makin keras zatnya, makin kasar pemaksaan simpangnya
        --- dan makin cepat perangkaian ulang ganjaran yang
        menjadi mesin \omterm{prop:g5:healthy-choices:substances}{kecanduannya}.
\end{itemize}
\end{proposition}
\admitted

\begin{remark}[Kecanduan, secara mekanis]\label{rem:g8:protecting-the-nervous-system:addiction}
Jebakan pada \cref{rem:g5:healthy-choices:never} kini punya
sebuah bagan kabel: simpang ganjarannya, yang ditakar
berulang-ulang, \emph{menyesuaikan diri} --- reseptornya
dipangkas, ambangnya dinaikkan --- sampai zatnya dibutuhkan
sekadar untuk merasa biasa, dan ketiadaannya menjadi protes
sistemnya. Penyesuaian itulah ketaksetangkupan
mudah-masuk-sukar-keluar yang terbuat dari \omterm{def:g8:nervous-communication:synapse}{sinapsis}:
memulainya adalah satu keputusan; berhentinya adalah melawan
simpang yang sudah dirangkai ulang. \omterm{def:g5:brain-in-command:system}{Otak} \omterm{def:g3:stages-of-life:stages}{remaja} yang masih
dapat disetel merangkai ulang paling cepat --- itulah
sebabnya setiap statistik menemukan pemula yang awal
terjebak paling dalam, dan sebabnya jangan-pernah-mulai
tetap menjadi langkah yang murah.
\end{remark}

\begin{example}[Gelung layarnya]\label{ex:g8:protecting-the-nervous-system:screens}
Tanpa molekul, simpang yang sama: umpan dan permainannya
disetel --- dengan sengaja --- untuk menggelitik kabel
ganjarannya dengan hadiah kecil yang tak terduga, jadwal
yang paling keras melatih simpangnya. Hasilnya adalah sepupu
ringan jebakannya: memeriksa yang sukar dihentikan, malam
yang lenyap, \omterm{prop:g1:healthy-habits:needs}{tidur} yang terdorong makin larut (jam pada
\cref{prop:g5:healthy-choices:screens}, ditambah ongkos pada
\cref{prop:g8:protecting-the-nervous-system:sleep}).
Pertahanannya sama seperti dahulu: putuskan jamnya lebih
dahulu, dan jagalah jam terakhirnya bebas layar.
\end{example}

\section{Praktik: menjalankan perlindunganmu sendiri}

\begin{method}[Buku panduan pemilik sistem sarafnya]\label{met:g8:protecting-the-nervous-system:manual}
Daftar perawatan pada \cref{met:g5:healthy-choices:run},
diterbitkan ulang beserta mekanismenya:
\begin{enumerate}
  \item \emph{tidurlah sekian jamnya} --- pemadatan dan tata
        graha berjalan pada malam hari; jagalah jam
        terakhirnya dari layar;
  \item \emph{anggarkan kenyaringannya} --- modal
        penerimanya tidak tumbuh kembali: kenyaringan
        diturunkan, jeda disisipkan, sumbat \omterm{def:g1:the-five-senses:sense}{telinga} dipakai
        di tempat yang menderu;
  \item \emph{helm} di tempat jatuhnya cepat
        (\cref{ex:g5:brain-in-command:helmet}) --- kabelnya
        tidak pulih seperti \omterm{def:g3:bones-and-muscles:skeleton}{tulang};
  \item \emph{tanpa pemerdaya simpang} bagi \omterm{def:g5:brain-in-command:system}{otak} yang sedang
        dibangun --- dan ``tidak''-nya diputuskan dengan
        tenang lebih dahulu (kalimat saku pada
        \cref{ex:g5:healthy-choices:answers}, kini dengan
        alasan lengkapnya);
  \item \emph{pakailah mesinnya} --- belajar, olahraga,
        musik: simpangnya menguat oleh lalu lintas, dan
        gelung yang terlatih adalah harta yang dilindungi
        setiap kebiasaan lainnya.
\end{enumerate}
\end{method}

\begin{example}[Pertolongan ada, dan berhasil]\label{ex:g8:protecting-the-nervous-system:help}
Bagi siapa pun yang sudah terjerat --- kebiasaannya sendiri
atau kebiasaan seorang teman, zat atau layar: alamat pada
\cref{ex:g8:contraception:help} melayani di sini juga.
Perawat sekolah dan dokternya memperlakukan \omterm{prop:g5:healthy-choices:substances}{kecanduan}
sebagaimana bab ini memperlihatkannya --- sebuah keadaan
kabel, bukan vonis atas watak --- dan simpang yang sudah
dirangkai ulang dapat menyesuaikan diri kembali: lebih sukar
daripada tidak pernah memulai, jauh dari mustahil. Rencana
yang paling buruk adalah rencana halaman sekolahnya: diam.
\end{example}

\section{Latihan}

\begin{exercise}[$\star$]\label{exo:g8:protecting-the-nervous-system:1}
Apa yang dikerjakan \omterm{prop:g1:healthy-habits:needs}{tidur} di simpangnya, dan siapa yang
paling membutuhkan jamnya?
\end{exercise}

\begin{exercise}[$\star$]\label{exo:g8:protecting-the-nervous-system:2}
Mengapa bunyi yang keras adalah soal anggaran? Apa yang
dibelanjakan, dan apa peringatan takaran berlebihnya?
\end{exercise}

\begin{exercise}[$\star$]\label{exo:g8:protecting-the-nervous-system:3}
Di mana semua zat psikoaktif bekerja, dan dengan tiga cara
yang mana?
\end{exercise}

\begin{exercise}[$\star$]\label{exo:g8:protecting-the-nervous-system:4}
Apa yang dilakukan alkohol terhadap gelungnya --- sebutkan
tiga perhentiannya.
\end{exercise}

\begin{exercise}[$\star$]\label{exo:g8:protecting-the-nervous-system:5}
Nyatakan mekanisme \omterm{prop:g5:healthy-choices:substances}{kecanduannya} dalam dua kalimat.
\end{exercise}

\begin{exercise}[$\star$]\label{exo:g8:protecting-the-nervous-system:6}
Mengapa gelung layarnya termasuk dalam bab ini, walaupun
tanpa molekul?
\end{exercise}

\begin{exercise}[$\star\star$]\label{exo:g8:protecting-the-nervous-system:7}
Jelaskan paradoks belajar ulangnya dengan
\cref{prop:g8:protecting-the-nervous-system:sleep}.
\end{exercise}

\begin{exercise}[$\star\star$]\label{exo:g8:protecting-the-nervous-system:8}
Mengapa takaran yang sama lebih mahal bagi \omterm{def:g5:brain-in-command:system}{otak} \omterm{def:g3:stages-of-life:stages}{remaja}
daripada bagi \omterm{def:g5:brain-in-command:system}{otak} dewasa? Dua mekanisme
(\cref{rem:g8:puberty:moods},
\cref{rem:g8:protecting-the-nervous-system:addiction}).
\end{exercise}

\begin{exercise}[$\star\star$]\label{exo:g8:protecting-the-nervous-system:9}
Hubungkan meter pada \cref{ex:g5:brain-in-command:driver}
dengan peredaman alkoholnya: apa yang memanjang, dan mengapa
undang-undang mabuk-mengemudi adalah undang-undang \omterm{def:g5:brain-in-command:reaction}{waktu reaksi}?
\end{exercise}

\begin{exercise}[$\star\star$]\label{exo:g8:protecting-the-nervous-system:10}
Berikan alasan ulang bagi tiga butir buku panduan pemiliknya
dengan mekanisme setingkat simpangnya.
\end{exercise}

\begin{exercise}[$\star\star$]\label{exo:g8:protecting-the-nervous-system:11}
``Sebuah keadaan kabel, bukan vonis atas watak.'' Apa yang
diubah pembingkaian ulang itu tentang mencari pertolongan
--- dan tentang menawarkannya kepada seorang teman?
\end{exercise}

\begin{exercise}[$\star\star\star$]\label{exo:g8:protecting-the-nervous-system:12}
Tuliskan hujah penutup tahunnya: dari celah yang dapat
disetel pada \cref{def:g8:nervous-communication:synapse},
turunkanlah sekaligus kekuatan terbesar \omterm{def:g5:brain-in-command:system}{otaknya} dan
kerawanan khasnya --- belajar dan \omterm{prop:g5:healthy-choices:substances}{kecanduan} sebagai satu
mekanisme yang dibaca dua kali.
\end{exercise}

\section{Soal: Malam Pencegahan}

\begin{problem}[{Soal akhir pekan --- malam orang
tua-dan-murid di sekolahnya, dinaskahkan beserta
mekanismenya}]\label{pb:g8:protecting-the-nervous-system:1}
Kelasnya menyiapkan malam pencegahan sekolahnya: empat
perhentian --- \omterm{prop:g1:healthy-habits:needs}{tidur}, bunyi, zat, layar --- masing-masing
dengan sebuah peragaan dan sebuah kartu mekanisme.

\medskip\noindent\textbf{Bagian I --- Kartu
perhentiannya.}
\begin{enumerate}
  \item Tuliskan kartu \omterm{prop:g1:healthy-habits:needs}{tidurnya}: peragaannya adalah nilai
        uji penggaris dua sukarelawan
        (\cref{met:g5:brain-in-command:ruler}), yang
        beristirahat cukup melawan yang bermalam pendek.
        Ramalkan nilainya lalu berikan mekanisme simpangnya.
  \item Tuliskan kartu bunyinya: sebuah pengukur desibel,
        sepasang sumbat \omterm{def:g1:the-five-senses:sense}{telinga}, dan kalimat ``kamu tidak
        dapat menumbuhkannya kembali''. Isilah mekanismenya.
  \item Tuliskan kartu zatnya: satu bagan sebuah \omterm{def:g8:nervous-communication:synapse}{sinapsis},
        tiga panah berlabel menirukan, menghadang, memaksa.
        Tetapkan alkohol, nikotin dan sebuah bisa dari
        \cref{exo:g8:nervous-communication:11} pada panahnya.
  \item Tuliskan kartu layarnya: tanpa molekul --- jadi apa
        yang dibagi bersama antara peragaan
        telepon-dengan-pemberitahuan pada perhentiannya dan
        perhentian zatnya?
\end{enumerate}

\medskip\noindent\textbf{Bagian II --- Pertanyaan dari
lantainya.}
\begin{enumerate}[resume]
  \item Seorang orang tua: ``kami minum pada usia mereka dan
        ternyata baik-baik saja.'' Jawablah dengan
        ketaksetangkupan lokasi pembangunannya --- apa yang
        berbeda pada usia empat belas?
  \item Seorang murid: ``satu batang rokok tak mungkin
        merangkai ulang apa pun.'' Jawablah dengan
        \cref{rem:g8:protecting-the-nervous-system:addiction}
        dan \cref{rem:g5:healthy-choices:never} --- apa harga
        sebenarnya takaran pertamanya?
  \item Seorang orang tua: ``bukankah layar hanyalah
        televisi yang baru?'' Bedakan jadwalnya: apa yang
        membuat hadiah kecil yang tak terduga menjadi
        pelatih yang lebih keras?
  \item Seorang murid, dengan pelan: ``dan bila seseorang
        sudah tak dapat berhenti?'' Berikan jawaban
        \cref{ex:g8:protecting-the-nervous-system:help},
        beserta alamatnya.
\end{enumerate}

\medskip\noindent\textbf{Bagian III --- Ceramah
penutupnya.}
\begin{enumerate}[resume]
  \item Ceramahnya dibuka dengan sang penyulap bola dan
        ditutup dengan \omterm{prop:g5:healthy-choices:substances}{kecanduan} --- ``keterstelan yang
        sama, dibaca dua kali''. Susunlah hujah itu dalam
        empat kalimat.
  \item Ia lalu menyerahkan kalimat saku pada
        \cref{ex:g5:healthy-choices:answers} kepada setiap
        murid. Jelaskan, beserta mekanismenya, mengapa
        kalimat itu harus disusun \emph{sekarang}.
  \item Poster malamnya memperlihatkan kelima butir buku
        panduan pemiliknya. Urutkan untuk sepekan seorang
        anak empat belas tahun, yang paling melindungi lebih
        dahulu, lalu pertahankan pilihan pertamamu.
  \item Tutuplah malamnya dalam satu kalimat: apa yang
        sedang dilindungi, dan mengapa empat belas adalah
        pertengahan dasawarsa yang menentukan?
\end{enumerate}
\end{problem}
